\documentclass[aspectratio=169,11pt]{beamer}

\usetheme{Madrid}
\usecolortheme{default}
\usefonttheme{professionalfonts}
\setbeamertemplate{navigation symbols}{}
\setbeamertemplate{footline}[frame number]

\usepackage{amsmath, amssymb, booktabs}

\title{Firms, Hiring, Management, And Organizational Response To New Technology}
\subtitle{AI Adoption As A Firm-Side Labor Problem}
\author{Special Topic 5: Technology, Innovation, and Labor -- Week 4}
\date{}

\begin{document}

\begin{frame}
  \titlepage
\end{frame}

\begin{frame}{Central question}
\begin{itemize}
    \item When a firm adopts a new technology, what changes inside the firm?
    \item Which jobs are redesigned, which workers are trained, and which vacancies are posted?
    \item How do supervision, monitoring, delegation, promotion, and internal mobility change?
    \item Week 4 discipline: adoption is not complete until work is reorganized.
\end{itemize}
\end{frame}

\begin{frame}{Why this is labor economics}
\begin{itemize}
    \item Technology is implemented through vacancies, task assignments, managers, teams, and contracts.
    \item The same AI tool can substitute for some tasks and augment others.
    \item Firm choices determine whether incumbents are retrained, replaced, monitored, or promoted.
    \item Worker and manager attitudes can shape usage, resistance, trust, and measured productivity.
\end{itemize}
\end{frame}

\begin{frame}{From adoption to reorganization}
\small
\begin{tabular}{p{0.22\linewidth} p{0.38\linewidth} p{0.25\linewidth}}
\toprule
Margin & What changes & Labor object \\
\midrule
Acquisition & buy, build, license, or hire for AI & investment and adoption \\
Implementation & connect data, workflows, incentives & management and training \\
Usage & workers or managers rely on the tool & task performance and quality \\
Reorganization & jobs, teams, authority, ladders change & labor demand and incidence \\
\bottomrule
\end{tabular}
\vspace{0.35cm}
\begin{itemize}
    \item Exposure, adoption, usage, and redesign are different empirical treatments.
\end{itemize}
\end{frame}

\begin{frame}{Adoption as an organizational choice}
\small
\[
Adopt_{ft}=1
\left[
E_t\left(\Pi_{ft}^{AI}(K_{ft}^{O}, H_{ft}, M_{ft}, A_{ft})\right)
- \Pi_{ft}^{0}
- C_{ft}^{adopt}
- C_{ft}^{org}
> 0
\right]
\]
\begin{itemize}
    \item $K^O$: organizational capital; $H$: workforce skills and hierarchy.
    \item $M$: management practice; $A$: worker and manager acceptance or resistance.
    \item Adoption depends on complementary investments, not only the tool price.
\end{itemize}
\end{frame}

\begin{frame}{Adopters are selected}
\begin{itemize}
    \item Early AI adopters may already be larger, more productive, more skilled, and better managed.
    \item Adoption estimates mix the technology with firm capability, timing, and workforce composition.
    \item A vacancy or resume measure captures intended skill demand, not necessarily realized usage.
    \item A survey adoption measure may capture aspiration, experimentation, or full workflow redesign.
\end{itemize}
\end{frame}

\begin{frame}{AI as three labor objects}
\small
\begin{tabular}{p{0.24\linewidth} p{0.37\linewidth} p{0.26\linewidth}}
\toprule
Margin & Mechanism & Labor implication \\
\midrule
Substitution & AI performs tasks cheaply & task displacement or role compression \\
Augmentation & AI raises worker productivity & higher output, quality, or scale \\
Delegation & AI changes who decides & authority, monitoring, and evaluation shift \\
\bottomrule
\end{tabular}
\vspace{0.35cm}
\begin{itemize}
    \item The same tool can substitute, augment, and reallocate control at the same firm.
\end{itemize}
\end{frame}

\begin{frame}{Firm labor-demand decomposition}
\[
\Delta L_f =
\Delta L_f^{sub}
+ \Delta L_f^{aug}
+ \Delta L_f^{scale}
+ \Delta L_f^{org}
\]
\begin{itemize}
    \item Substitution: fewer workers needed for automated tasks.
    \item Augmentation: worker-performed tasks become more productive.
    \item Scale: lower costs or higher quality raise output demand.
    \item Organization: hierarchy, monitoring, teams, and job boundaries change.
\end{itemize}
\end{frame}

\begin{frame}{Job redesign and task reassignment}
\begin{itemize}
    \item Firms rebundle tasks into redesigned jobs rather than replacing occupations one for one.
    \item Workers may shift toward validation, exception handling, customer communication, or system maintenance.
    \item The key question is whether old firm-specific experience remains valuable after redesign.
    \item Measured wage gains can reflect recomposition toward new workers rather than incumbent gains.
\end{itemize}
\end{frame}

\begin{frame}{Hiring, screening, and internal labor markets}
\begin{itemize}
    \item Hiring response: AI engineers, data roles, compliance roles, domain experts, and AI-fluent generalists.
    \item Screening response: more emphasis on autonomy, judgment, digital fluency, and auditing machine output.
    \item Internal response: training, task reassignment, lateral moves, promotion redesign, or exit.
    \item Incumbent training and hiring around incumbents are different labor-market mechanisms.
\end{itemize}
\end{frame}

\begin{frame}{Principal-agent problems}
\small
\begin{tabular}{p{0.24\linewidth} p{0.31\linewidth} p{0.31\linewidth}}
\toprule
Margin & Firm or manager objective & Worker concern \\
\midrule
Monitoring & observe effort and output & privacy, autonomy, gaming \\
Evaluation & standardize performance metrics & fairness and opacity \\
Delegation & speed and consistency & loss of discretion \\
Replacement risk & lower costs and reassign tasks & job security and de-skilling \\
\bottomrule
\end{tabular}
\vspace{0.35cm}
\begin{itemize}
    \item AI changes not only productivity, but who controls information and who bears risk.
\end{itemize}
\end{frame}

\begin{frame}{AI attitudes and algorithm aversion}
\begin{itemize}
    \item Workers may resist AI because of replacement risk, monitoring, unfair evaluation, or loss of control.
    \item Managers may resist delegation because of overconfidence, liability, discretion, or status.
    \item Attitudes affect measured adoption: access, self-reported use, mandated use, and trusted use differ.
    \item Algorithm aversion is a labor issue when it changes hiring, screening, promotion, or supervision.
\end{itemize}
\end{frame}

\begin{frame}{Frontier methods}
\small
\begin{tabular}{p{0.31\linewidth} p{0.40\linewidth} p{0.17\linewidth}}
\toprule
Design & Identifies well & Main caution \\
\midrule
Worker survey + outcomes & trust, fear, usage, retention & attitudes are endogenous \\
Management survey + admin data & beliefs and implementation strategy & manager quality confounds \\
Randomized rollout & access or encouragement effects & short-run horizon \\
Vignette or conjoint & fairness, transparency, control & stated preference \\
Usage logs + surveys & realized use and acceptance & compliance vs belief \\
Employer-worker panel & worker incidence around adoption & adopter selection \\
\bottomrule
\end{tabular}
\end{frame}

\begin{frame}{What a strong design separates}
\begin{enumerate}
    \item Technology availability or exposure.
    \item Firm adoption and implementation timing.
    \item Actual worker or manager usage.
    \item Attitudes: trust, fairness, replacement concern, delegation willingness.
    \item Organizational redesign: hiring, training, monitoring, promotion, and job boundaries.
    \item Worker outcomes: productivity, wages, separations, autonomy, job quality, and mobility.
\end{enumerate}
\end{frame}

\begin{frame}{Week 4 lab design}
\begin{itemize}
    \item Reproduce: build a synthetic firm AI-adoption panel with workforce composition and hierarchy.
    \item Diagnose: separate selection, acquisition, usage, substitution, augmentation, scale, and redesign.
    \item Transfer: add worker and manager attitudes inspired by algorithm-aversion designs.
    \item Output: compact tables that distinguish adoption, usage, sentiment, productivity, and worker incidence.
    \item Goal: practice the design logic without confidential firm records or official replication data.
\end{itemize}
\end{frame}

\begin{frame}{Bridge to Week 5}
\begin{itemize}
    \item Week 4 studies how firms implement technology inside jobs and internal labor markets.
    \item Week 5 asks how firm-level adoption aggregates into inequality, rents, market power, and institutions.
    \item The bridge is distributional: firm capabilities and worker access shape who captures technology gains.
\end{itemize}
\end{frame}

\begin{frame}{Takeaways}
\begin{itemize}
    \item Firm adoption is an organizational labor problem, not a simple capital purchase.
    \item AI can substitute, augment, and reallocate authority at the same time.
    \item Hiring, training, monitoring, promotion, and job redesign determine worker incidence.
    \item Worker and manager attitudes are adoption frictions and inequality channels.
    \item Strong empirical work separates adoption, usage, sentiment, redesign, and outcomes.
\end{itemize}
\end{frame}

\end{document}
